\abstract{
  Violations of Bell inequalities are commonly interpreted as evidence
  for nonlocal influences or as constraints on realist descriptions.
  We show that the failure of Bell-type factorizability arises generically
  whenever observable outcomes are obtained through a non-injective mapping
  from an underlying configuration space.
  In this setting, the standard factorization assumption can be viewed as
  an implicit requirement that observable variables admit a jointly
  factorizable completion at the underlying level.
  We demonstrate that this requirement fails structurally when the mapping
  from underlying configurations to observables is many-to-one.
  The resulting breakdown of probabilistic factorization does not rely on
  superluminal dynamics or hidden causal influences, but follows from
  information loss under projection.
  Observable outcomes correspond to equivalence classes of underlying
  configurations, preventing the assignment of independent local variables.
  We illustrate this mechanism with an explicit toy model producing
  Bell--CHSH violations while preserving operational no-signalling and
  statistical independence of measurement settings.
  The model is not intended to reproduce quantum correlations quantitatively,
  and may exceed the Tsirelson bound; its role is to isolate the structural
  origin of the violation.
  The analysis preserves locality at the underlying level, introduces no
  additional hidden-variable dynamics, and does not modify quantum mechanics.
  It identifies non-injective projection as a sufficient structural condition
  for the failure of Bell factorizability and clarifies how classical
  factorization is recovered in regimes where the effective mapping becomes
  approximately injective.
}
